\section*{Resumen.}

El proyecto Patitas Felices consiste en el desarrollo de un sistema web para la gestión y administración de una clínica veterinaria, cuyo propósito es mejorar la organización de sus actividades y facilitar el manejo de la información relacionada con los propietarios y sus mascotas. El sistema busca reducir los procesos manuales, optimizar el tiempo de atención y mantener un registro organizado y accesible de los servicios veterinarios.

La plataforma contará con seis módulos principales: propietarios y mascotas, agenda y disponibilidad, atenciones y expediente clínico, vacunación y seguimientos, insumos e inventario, y portal y reportes. Además, contará con diferentes tipos de usuarios: administrador, recepcionista, veterinario y propietario, cada uno con funciones y permisos específicos. Para su desarrollo se utilizan tecnologías como PHP, Laravel, SQLite, HTML, CSS y JavaScript, que permiten construir una aplicación web organizada, funcional y adaptable a las necesidades de la clínica.

En conclusión, Patitas Felices busca ofrecer una solución tecnológica que facilite la administración de una clínica veterinaria, mejore el control de la información y contribuya a brindar una atención más organizada y eficiente a las mascotas y sus propietarios.

